%%%%%%%%%%%%%%%%%%%%%%%%%%%%%%%%%%%%%%%%%%%%%%%%%%%%
\usepackage[english]{babel}
\usepackage[utf8]{inputenc}
\usepackage[T1]{fontenc}
\usepackage{lmodern}

\usepackage{hyperref} % must be loaded before glossaries-extra

% bibliography
\usepackage[hyperref=true,backend=biber,maxbibnames=9,maxcitenames=2,style=numeric-comp,sorting=ynt,backref=false]{biblatex} % hyperref uses links, backref goes back to citations, uses biber as backend, with 9 names at most in bibliography and 2 in citations, citing using numbers, and sorting in citation order
% sorting can be also ydnt for year descending, name, title or ynt for ascending year
% Rimuove corsivo e virgolette da tutti i titoli principali
\DeclareFieldFormat{title}{#1}
\DeclareFieldFormat[inbook]{title}{#1}       % necessario separatamente per inbook
\DeclareFieldFormat{booktitle}{#1}
\DeclareFieldFormat{journaltitle}{#1}
\DeclareFieldFormat[article]{title}{#1}
\DeclareFieldFormat[article]{year}{#1}
\DeclareFieldFormat[incollection]{title}{#1}
\DeclareFieldFormat[inproceedings]{title}{#1}
% Rimuove le parentesi dall'anno per gli articoli
\renewbibmacro*{issue+date}{%
  \setunit*{\addcomma\space}%
  \printfield{year}%
  \newunit}

% Cambia separatore dei campi da punto a virgola]
\renewcommand*{\newunitpunct}{\addcomma\space}

\usepackage{adjustbox} % to resize boxes by keeping the same aspect ratio
\usepackage{algorithm} % algorithm environment
\usepackage{algpseudocode} % improved pseudo-code
\usepackage{amsfonts}               %  AMS mathematical fonts
\usepackage{amsmath}
\usepackage{amssymb}                %  AMS mathematical symbols
\usepackage{bm}                     %  black/bold mathematical symbols
\usepackage{booktabs}               %  better tables
\usepackage[labelfont=bf]{caption} % font=footnotesize % to have reduced caption font size
\usepackage{csquotes}
\usepackage{enumitem} %left align the bulleted points
\usepackage{geometry}
%\usepackage{glossaries} % to use acronyms and glossary, it has also glossaries-extra as extension, but commands are different
\usepackage[%
    toc, % puts the link in the ToC
    %record, % to use bib2gls
    abbreviations, % to load abbreviations / acronyms
    nonumberlist, % to avoid printing the numbers of the references in the acronyms page
]{glossaries-extra}
\usepackage{graphicx}               %  post-script images
%\usepackage{iwona} % extra fonts, substitute standard ones
\usepackage{listings} % to insert formatted code
\usepackage{lipsum} % for lorem ipsum text, not needed in the real work
\usepackage{makecell} % to change dimensions of cells, for math cases
\usepackage{mathtools} % for additional commands
\usepackage{mfirstuc} % to have capitalization capabilities
\usepackage[final]{microtype}      % microtypography, final lets latex use it also in bibliography
\usepackage{multirow} % to allow for cells covering more than 1 row in tables
\usepackage{nicefrac}       % compact symbols for 1/2, etc.
%\usepackage[lofdepth,lotdepth]{subfig}
\usepackage{ragged2e} % for justifying text
\usepackage{siunitx} % support for SI units of measurement and number typesetting
\usepackage{subfig}
\usepackage{svg} % for svg support, works only if inkscape is installed, default for Overleaf v2
%\usepackage{subfigure}              %  subfigure compatibility, can be removed if subfig
\usepackage{tabularx} % equal-width columns in tables
\usepackage{textcomp} % extra fonts and symbols
\usepackage{url}            % simple URL typesetting
\usepackage{verbatim} % for extended verbatim support
\usepackage{xcolor} % to define colors and use standard CSS names add dvipsnames as option, but it clashes with xcolor loaded in toptesi, pay attention that if it goes in conflict with tikz/beamer, simply use \documentclass[usenames,dvipsnames]{beamer}, along with other custom options when defining the document class

\usepackage{tikz-cd} % diagrammi commutativi
\usepackage{tikz}
\usepackage{amsthm}
\usepackage{hyperref}
\usepackage[most]{tcolorbox}
\usepackage[capitalize]{cleveref} % Per riferimenti automatici e corretti
\usepackage{fancyhdr}

% Stile di base per tcolorbox
\tcbset{
  base/.style={
    arc=0mm, 
    bottomtitle=0.5mm,
    boxrule=0mm,
    colbacktitle=black!10!white, 
    coltitle=black, 
    fonttitle=\bfseries, 
    left=2.5mm,
    leftrule=1mm,
    right=3.5mm,
    title={#1},
    toptitle=0.75mm,
  }
}


% Box per teoremi (Theorems, Lemmas, etc.)
\newtcolorbox{mainbox}[1]{
  enhanced,
  parbox=false,
  sharp corners,
  colback=gray!10!white,        % Sfondo grigio chiaro
  colframe=white,               % Nessun bordo esterno
  boxrule=0pt,                  % Rimuove il bordo attorno
  borderline west={2.5pt}{0mm}{black!65!white}, % Linea a sinistra dello stesso grigio
  left=3.5mm,                   % Spazio a sinistra
  right=3.5mm,                  % Spazio a destra
  breakable,
  title={#1}
}

\newtcolorbox{subbox}[1]{
  enhanced,
  parbox=false,
  sharp corners,
  colback=white,             % Sfondo bianco
  colframe=black!30!white,   % Colore usato solo per leftrule
  boxrule=0pt,               % Nessuna cornice visibile
  frame hidden,              % Nasconde il frame completamente
  borderline west={2.5pt}{0pt}{black!10!white}, % Riga verticale a sinistra
  breakable,
  left=3.5mm,                   % Spazio a sinistra
  right=3.5mm, 
  title={#1}                 % Titolo da parametro
}

% Remarks box
\newtcolorbox{rembox}[1]{
  enhanced,
  parbox=false,
  sharp corners,
  colback=white,             % Sfondo bianco
  colframe=black!30!white,   % Colore usato solo per leftrule
  boxrule=0pt,               % Nessuna cornice visibile
  frame hidden,              % Nasconde il frame completamente
  borderline west={2.5pt}{0pt}{black!40!white}, % Riga verticale a sinistra
  breakable,
  left=3.5mm,                   % Spazio a sinistra
  right=3.5mm,
  title={#1}                 % Titolo da parametro
}

% Definizione degli ambienti per Teoremi, Lemmi, Corollari, etc.
\newtheorem{theo}{Theorem}[section]
\newtheorem{lemma}[theo]{Lemma}
\newtheorem{cor}[theo]{Corollary}
\newtheorem{prop}[theo]{Proposition}
\newtheorem{propr}[theo]{Property}
\newtheorem{defi}[theo]{Definition}
\newtheorem{rem}[theo]{Remark}
\newtheorem{ex}[theo]{Example}

% Etichette per i riferimenti di cleveref
\crefname{theo}{Theorem}{Theorems}
\crefname{lemma}{Lemma}{Lemmas}
\crefname{cor}{Corollary}{Corollaries}
\crefname{prop}{Proposition}{Propositions}
\crefname{propr}{Property}{Properties}
\crefname{defi}{Definition}{Definitions}
\crefname{rem}{Remark}{Remarks}
\crefname{ex}{Example}{Examples}

\Crefname{theo}{Theorem}{Theorems}
\Crefname{lemma}{Lemma}{Lemmas}
\Crefname{cor}{Corollary}{Corollaries}
\Crefname{prop}{Proposition}{Propositions}
\Crefname{propr}{Property}{Properties}
\Crefname{defi}{Definition}{Definitions}
\Crefname{rem}{Remark}{Remarks}
\Crefname{ex}{Example}{Examples}

\Crefformat{theo}{\textit{Theorem~#2#1#3}}
\Crefformat{lemma}{\textit{Lemma~#2#1#3}}
\Crefformat{cor}{\textit{Corollary~#2#1#3}}
\Crefformat{prop}{\textit{Proposition~#2#1#3}}
\Crefformat{propr}{\textit{Property~#2#1#3}}
\Crefformat{defi}{\textit{Definition~#2#1#3}}
\Crefformat{rem}{\textit{Remark~#2#1#3}}
\Crefformat{ex}{\textit{Example~#2#1#3}}


% Cancella intestazioni e piè di pagina
\fancyhf{}


% Linea sopra l'intestazione
\renewcommand{\headrulewidth}{0.4pt}
\renewcommand{\footrulewidth}{0pt}